\chapter{2012年广东卷（理）}

\section{填空题}

\begin{problem}
不等式 \(|x + 2| - |x| \leqslant 1\) 的解集为 \fillinblank{} \fillinblank{}.
\end{problem}

\begin{problem}
\(\left(x^{2} + \frac{1}{x}\right)^{6}\) 的展开式中 \(x^{3}\) 的系数为 \fillinblank{}. (用数字作答)
\end{problem}

\begin{problem}
已知递增的等差数列 \(\{a_{n}\}\) 满足 \(a_{1} = 1, a_{3} = a_{2}^{2} - 4\)，则 \(a_{n} = \)  \fillinblank{}.
\end{problem}

\begin{problem}
曲线 \( y = x^{3} - x + 3 \) 在点(1,3)处的切线方程为 \fillinblank{} \fillinblank{}.
\end{problem}

\begin{problem}
执行如图所示的程序框图，若输入 \(n\) 的值为 8，则输出 \(s\) 的值为 \fillinblank{}.

\bitmapfigure[max width=0.38\linewidth,max totalheight=6cm]{img/2012/guangdong_science/q13_fig1.png}
\end{problem}

\begin{problem}
在平面直角坐标系 \(xOy\) 中，曲线 \(C_1\) 和 \(C_2\) 的参数方程分别为 \(\left\{ \begin{array}{l} x = t, \\ y = \sqrt{t}, \end{array} \right.\) (t 为参数) 和 \(\left\{ \begin{array}{l} x = \sqrt{2}\cos \theta, \\ y = \sqrt{2}\sin \theta, \end{array} \right.\) (\(\theta\) 为参数)，则曲线 \(C_1\) 和 \(C_2\) 的交点坐标为 \fillinblank{}.
\end{problem}

\begin{problem}
如图，圆 \(O\) 的半径为 1, \(A, B, C\) 是圆周上的三点. 满足 \(\angle ABC = 30^{\circ}\)，过点 \(A\) 作圆 \(O\) 的切线与 \(OC\) 的延长线交于点 \(P\)，则 \(PA = \fillinblank{}.\)

\bitmapfigure[max width=0.45\linewidth,max totalheight=5.8cm]{img/2012/guangdong_science/q15_fig1.png}
\end{problem}

\section{解答题}

\begin{problem}
已知函数 \( f(x) = 2\cos \left( \omega x + \frac{\pi}{6} \right) \) (其中 \( \omega > 0, x \in \R \) 的最小正周期为 \( 10\pi \)). (1) 求 \( \omega \) 的值; (2) 设 \( \alpha, \beta \in \left[0, \frac{\pi}{2}\right] \)，\( f\left( 5\alpha + \frac{5}{3}\pi \right) = -\frac{6}{5} \)，\( f\left( 5\beta - \frac{5}{6}\pi \right) = \frac{16}{17} \)，求 \( \cos (\alpha + \beta) \) 的值.
\end{problem}

\begin{problem}
\bitmapfigure[max width=0.42\linewidth,max totalheight=5.8cm]{img/2012/guangdong_science/q17_fig1.png}

某班 50 位学生期中考试数学成绩的频率分布直方图如图所示，其中成绩分组区间是: [40,50)，[50,60), [60,70), [70,80), [80,90), [90,100].

(1) 求图中 \(x\) 的值; (2) 从成绩不低于 80 分的学生中随机选取 2 人，该 2 人中成绩在 90 分以上 (含 90 分) 的人数记为 \(\xi\)，求 \(\xi\) 的数学期望.
\end{problem}

\begin{problem}
\bitmapfigure[max width=0.46\linewidth,max totalheight=5.8cm]{img/2012/guangdong_science/q18_fig1.png}

如图所示，在四棱锥 \(P - ABCD\) 中，底面 \(ABCD\) 为矩形，\(PA \perp\) 平面 \(ABCD\)，点 \(E\) 在线段 \(PC\) 上，\(PC \perp\) 平面 \(BDE\). (1) 证明: \(BD \perp\) 平面 \(PAC\); (2) 若 \(PA = 1, AD = 2\)，求二面角 \(B - PC - A\) 的正切值.
\end{problem}

\begin{problem}
在平面直角坐标系 \(xOy\) 中，已知椭圆 \(C: \frac{x^2}{a^2} + \frac{y^2}{b^2} = 1 (a > b > 0)\) 的离心率 \(e = \sqrt{\frac{2}{3}}\)，且椭圆 \(C\) 上的点到点 \(Q(0,2)\) 的距离的最大值为 3. (1) 求椭圆 C 的方程; (2) 在椭圆 \(C\) 上，是否存在点 \(M(m, n)\)，使得直线 \(l: mx + ny = 1\) 与圆 \(O: x^2 + y^2 = 1\) 相交于不同的两点 \(A, B\)，且 \(\triangle OAB\) 的面积最大？若存在，求出点 \(M\) 的坐标及相对应的 \(\triangle OAB\) 的面积; 若不存在，请说明理由.
\end{problem}

\begin{problem}
设 \(a < 1\)，集合 \(A = \{x \in \R \mid x > 0\}\)，\(B = \{x \in \R \mid 2x^2 - 3(1 + a)x + 6a > 0\}\)，\(D = A \cap B\). (1) 求集合 \(D\) (用区间表示); (2) 求函数 \( f(x) = 2x^3 - 3(1 + a)x^2 + 6ax \) 在 \( D \) 内的极值点.
\end{problem}

\begin{problem}
设数列 \(\{a_{n}\}\) 的前 \(n\) 项和为 \(S_{n}\)，满足 \(2S_{n} = a_{n + 1} - 2^{n + 1} + 1, n \in \mathbf{N}^{*}\)，且 \(a_{1}, a_{2} + 5, a_{3}\) 成等差数列. (1) 求 \( a_{1} \) 的值; (2) 求数列 \(\{a_{n}\}\) 的通项公式; (3) 证明: 对一切正整数 \(n\)，有 \(\frac{1}{a_1} + \frac{1}{a_2} + \frac{1}{a_3} + \cdots + \frac{1}{a_n} < \frac{3}{2}\).
\end{problem}
